\documentclass[11pt,letterpaper]{article}

\usepackage[letterpaper,top=0.45in,bottom=0.5in,left=0.5in,right=0.5in]{geometry}
\usepackage{graphicx}
\usepackage{xcolor}
\usepackage{listings}
\usepackage{booktabs}
\usepackage{hyperref}
\usepackage{caption}
\usepackage{float}
\usepackage{titlesec}
\usepackage{amsmath}

\hypersetup{
    colorlinks=true,
    linkcolor=blue!60!black,
    urlcolor=blue!60!black,
}

\setlength{\parskip}{2pt}
\setlength{\parindent}{0pt}
\titlespacing*{\section}{0pt}{6pt}{2pt}
\titlespacing*{\subsection}{0pt}{4pt}{2pt}
\captionsetup{skip=2pt,font=small}

\graphicspath{{figures/}}

% -----------------------------------------------------------------------------
% Optional figures: if the screenshot has not yet been captured from
% Vivado/Vitis, fall back to a labelled framed box so the rest of the
% report still compiles cleanly.
% -----------------------------------------------------------------------------
\usepackage{tikz}
\newcommand{\optionalfig}[3][\textwidth]{%
    \IfFileExists{figures/#2}{%
        \includegraphics[width=#1]{#2}%
    }{%
        \begin{tikzpicture}
            \draw[dashed,thick] (0,0) rectangle (#1,4cm);
            \node[align=center] at (0.5*#1,2cm) {%
                \textit{placeholder:} \texttt{#2}\\[2pt]
                {\footnotesize #3}%
            };
        \end{tikzpicture}%
    }%
}

% -----------------------------------------------------------------------------
% Code listing styles
% -----------------------------------------------------------------------------
\definecolor{codebg}{RGB}{248,248,248}
\definecolor{codekw}{RGB}{0,0,180}
\definecolor{codecmt}{RGB}{0,120,0}
\definecolor{codestr}{RGB}{170,55,20}
\definecolor{codenum}{RGB}{120,120,120}

\lstset{
    language=C,
    basicstyle=\ttfamily\footnotesize,
    keywordstyle=\color{codekw}\bfseries,
    commentstyle=\color{codecmt}\itshape,
    stringstyle=\color{codestr},
    numbers=left,
    numberstyle=\tiny\color{codenum},
    stepnumber=1,
    numbersep=5pt,
    backgroundcolor=\color{codebg},
    frame=single,
    rulecolor=\color{black!30},
    tabsize=4,
    showstringspaces=false,
    breaklines=true,
    breakatwhitespace=false,
    breakindent=0pt,
    postbreak=\mbox{\textcolor{red}{$\hookrightarrow$}\space},
    captionpos=b,
    aboveskip=4pt,
    belowskip=4pt,
    xleftmargin=6pt,
    xrightmargin=2pt,
    columns=fullflexible,
    keepspaces=true,
    morekeywords={u32,s32,Xil_Out32,Xil_In32}
}

\lstdefinelanguage{VHDLsmall}{
    sensitive=false,
    morekeywords={architecture,begin,case,component,downto,else,elsif,end,
        entity,if,is,library,others,port,map,signal,std_logic,
        std_logic_vector,then,to,use,when,work,all,generic,return,process,
        rising_edge,unsigned,signed,std_logic_unsigned,numeric_std,
        constant,variable,type,array,natural,integer,for,loop,while,resize},
    morecomment=[l]--
}

\lstdefinestyle{vhdlstyle}{
    language=VHDLsmall,
    basicstyle=\ttfamily\footnotesize,
    keywordstyle=\color{codekw}\bfseries,
    commentstyle=\color{codecmt}\itshape,
    numbers=left,
    numberstyle=\tiny\color{codenum},
    numbersep=5pt,
    backgroundcolor=\color{codebg},
    frame=single,
    rulecolor=\color{black!30},
    tabsize=4,
    breaklines=true,
    captionpos=b,
    aboveskip=4pt,
    belowskip=4pt,
    xleftmargin=6pt,
    xrightmargin=2pt,
    columns=fullflexible,
    keepspaces=true
}

\title{\vspace{-0.6in}\textbf{ECEC 661 -- Final}\\
       \large Stack Processor 101 with new \texttt{ssq} (``stack the square'') instruction}
\author{Gary Pham \texttt{(gp492)}}
\date{\today}

\begin{document}
\maketitle
\vspace{-0.25in}

% =============================================================================
\section{Problem Statement}
% =============================================================================

Starting from the assignment PDF's Stack Processor~101 (\texttt{sp101}),
implement a new machine instruction \texttt{ssq} -- \emph{stack the
square} -- which squares the value on top of the data stack in place.
The user-visible calling convention is:
\[
\texttt{sc}\ x\quad\texttt{ssq}\quad\texttt{halt}
\]
For \(x=7\) the processor must leave \(49\) on the stack and the stack
pointer must point to the next available address.  Because
\texttt{idle} initialises \(\mathit{sp}=128\) and \texttt{sc 7} writes
\(\mathrm{mem}[128]=7\) then bumps \(\mathit{sp}\) to~\(129\),
\texttt{ssq} overwrites \(\mathrm{mem}[128]\) with \(7\cdot7=49\) and
leaves \(\mathit{sp}=129\).  Deliverables, scored from
\href{run:../prob.md}{\texttt{Final/real/prob.md}}:

\begin{table}[H]
    \centering
    \small
    \begin{tabular}{lll}
        \toprule
        \textbf{Part} & \textbf{Points} & \textbf{Artifact in this report} \\
        \midrule
        1. simulate correctness (sim wave)           & 10 & Figs.~\ref{fig:sim1}--\ref{fig:simres2} \\
        2. custom IP user logic (component + port map) & 5 & Listing~\ref{lst:userlogic} \\
        3.1 block diagram                            & 3 & Fig.~\ref{fig:bd} \\
        3.2 C test app                               & 3 & Listing~\ref{lst:cmain} \\
        3.3 serial terminal snip (\(\mathrm{addr}\,128=\mathrm{square}\)) & 4 & Fig.~\ref{fig:vitis} / Listing~\ref{lst:uart} \\
        \bottomrule
    \end{tabular}
\end{table}

% =============================================================================
\section{Architecture}
% =============================================================================

The IP keeps the assignment PDF's three-block layout
(Figure~\ref{fig:bd}): an AXI4-Lite slave exposes the five
software-visible registers, a \texttt{bus\_ip\_mem\_bridge} multiplexes
the bus or the processor onto the single port of
\texttt{blk\_mem\_gen\_0}, and the stack processor (\texttt{user\_logic})
fetches/executes machine instructions out of that BRAM.  Only the
\texttt{exe} decoder of \texttt{user\_logic} is extended with the new
\texttt{ssq} opcode family; every other PDF instruction is unchanged.

\begin{figure}[H]
    \centering
    \optionalfig[0.92\textwidth]{diagram.png}{Vivado block design
        (Zynq PS, AXI Interconnect, \texttt{sp101\_axi\_0}).}
    \caption{Vivado block design: ZYNQ7 Processing System,
             Processor System Reset, AXI Interconnect, and the custom
             \texttt{sp101\_axi\_0} IP (\textit{Stack Processor 101 with
             ssq, AXI4-Lite}).  \texttt{M\_AXI\_GP0} from the PS reaches
             the IP through the interconnect; \texttt{FCLK\_CLK0}
             (100~MHz) and \texttt{FCLK\_RESET0\_N} drive the reset
             block.  The IP is mapped at base
             \texttt{0x43C3\_0000} (\(4\,\text{KB}\)) in the Address
             Editor.}
    \label{fig:bd}
\end{figure}

% =============================================================================
\section{Software Register Map}
% =============================================================================

The slave maps five 32-bit registers into a 32-byte aperture, identical
to the PDF's reference C snippet:

\begin{table}[H]
    \centering
    \caption{Software register map (byte offsets from
             \texttt{XPAR\_SP101\_AXI\_0\_S00\_AXI\_BASEADDR}).}
    \label{tab:regs}
    \small
    \begin{tabular}{clcp{8.0cm}}
        \toprule
        \textbf{Offset} & \textbf{Name} & \textbf{R/W} & \textbf{Purpose} \\
        \midrule
        \texttt{0x00} & \texttt{slv\_reg0} & R/W &
            control: bit~0~=\texttt{run}, bit~1~=\texttt{reset},
            bit~2~=\texttt{bus2mem\_en}, bit~3~=\texttt{bus2mem\_we} \\
        \texttt{0x04} & \texttt{slv\_reg1} & R/W &
            \texttt{bus2mem\_addr} (lower 10 bits) \\
        \texttt{0x08} & \texttt{slv\_reg2} & R/W &
            \texttt{bus2mem\_data\_in} (32 bit) \\
        \texttt{0x0C} & \texttt{slv\_reg3} & R   &
            \texttt{sp2bus\_data\_out} (32 bit) \\
        \texttt{0x10} & \texttt{slv\_reg4} & R   &
            bit~0 = \texttt{done} flag \\
        \bottomrule
    \end{tabular}
\end{table}

% =============================================================================
\section{Instruction Set and \texttt{ssq} Opcode}
% =============================================================================

The processor honours every opcode from the assignment PDF and adds a
new family for \texttt{ssq}.  The lowest hex digit of each constant is
the micro-step index inside the instruction; the higher hex digits are
the opcode group.

\begin{table}[H]
    \centering
    \caption{Instruction set (opcode entry constants in
             \href{run:../stackprocessor_101/rtl/user_logic.vhd}{\texttt{user\_logic.vhd}}).}
    \label{tab:isa}
    \small
    \begin{tabular}{llll}
        \toprule
        \textbf{Mnemonic} & \textbf{Code} & \textbf{Stack effect} & \textbf{Notes} \\
        \midrule
        \texttt{halt} & \texttt{0x000000FF} & --                                            & raise \texttt{done}, halt \\
        \texttt{sc}   & \texttt{0x00000001} & push \(c\)                                    & push constant pointed by pc \\
        \texttt{sl}   & \texttt{0x00000011} & pop \(a\), push \(\mathrm{mem}[a]\)           & indirect load \\
        \texttt{ss}   & \texttt{0x00000021} & pop \(d\), pop \(a\), \(\mathrm{mem}[a]\!\leftarrow\!d\) & indirect store \\
        \texttt{sadd} & \texttt{0x00000031} & pop, pop, push sum                            & stack add \\
        \texttt{ssq}  & \texttt{0x00000041} & pop \(x\), push \(x\cdot x\)                  & \textbf{NEW} -- unary in place \\
        \texttt{scp}  & \texttt{0x00000101} & pop \(s\), pop \(d\), \(\mathrm{mem}[d]\!\leftarrow\!\mathrm{mem}[s]\) & single-word copy \\
        \bottomrule
    \end{tabular}
\end{table}

The opcode \texttt{0x00000041} drops into the unused \texttt{0x4x}
slot above \texttt{sadd}, so none of the existing decoders are touched.
Because \texttt{ssq} is unary and reuses the same stack slot for input
and output, the stack pointer is unchanged: one ``pop'' and one
``push'' cancel.

\subsection{\texttt{ssq} micro-program}

The FSM walks \texttt{ssq} through six micro-states.  States
\texttt{SSQ}--\texttt{SSQ4} schedule the BRAM read of
\(\mathrm{mem}[\mathit{sp}-1]\) and burn the 4-cycle read latency
(BRAM~+~\texttt{mem\_data\_out\_pre}~+~\texttt{mem\_data\_out}).  State
\texttt{SSQ5} captures \(x\) in \texttt{temp1}; state \texttt{SSQ6}
issues the write of \(x\cdot x\) back to the same address and returns
to \texttt{fetch}.

\begin{lstlisting}[style=vhdlstyle,caption={Core micro-states (excerpt
of \texttt{when SSQ ...} body in
\href{run:../stackprocessor_101/rtl/user_logic.vhd}{\texttt{user\_logic.vhd}}).},label={lst:ssq}]
when SSQ =>                              -- start read of mem[sp-1]
    mem_addr <= std_logic_vector(unsigned(sp) - 1);
    we       <= "0"; ir <= SSQ2;
when SSQ2 => we <= "0"; ir <= SSQ3;      -- BRAM read latency
when SSQ3 => we <= "0"; ir <= SSQ4;      -- BRAM read latency
when SSQ4 => we <= "0"; ir <= SSQ5;      -- BRAM read latency
when SSQ5 =>                              -- x is now valid
    temp1 <= mem_data_out;
    we    <= "0"; ir <= SSQ6;
when SSQ6 =>                              -- write x*x back to mem[sp-1]
    mem_addr    <= std_logic_vector(unsigned(sp) - 1);
    mem_data_in <= std_logic_vector(
                       resize(unsigned(temp1) *
                              unsigned(temp1), 32));
    we          <= "1";
    n_s         <= fetch;
\end{lstlisting}

Multiplication uses \texttt{numeric\_std} (\(\texttt{unsigned}\times
\texttt{unsigned}\)) and \texttt{resize(\dots,32)} truncates the
64-bit product to the low 32 bits.  All test vectors used here have
\(x\le 46340\), so \(x\cdot x<2^{31}\) and no truncation occurs.

% =============================================================================
\section{Part 2 -- User Logic Instantiation (5 pts)}
% =============================================================================

\begin{sloppypar}
The required component declaration and port map live inside the
AXI4-Lite slave wrapper
\href{run:../stackprocessor_101/rtl/sp101_axi_v1_0_S00_AXI.vhd}%
{\texttt{sp101\_axi\_v1\_0\_S00\_AXI.vhd}}
(lines \(\sim 101\)--\(151\)).
Control bits are sliced straight from \texttt{slv\_reg0};
\texttt{bus2mem\_addr} takes the lower 10 bits of \texttt{slv\_reg1};
\texttt{bus2mem\_data\_in} is the full 32 bits of \texttt{slv\_reg2};
\texttt{sp2bus\_data\_out} and \texttt{done} feed the read mux
(\texttt{slv\_reg3} and \texttt{slv\_reg4}).
\end{sloppypar}

\begin{lstlisting}[style=vhdlstyle,caption={Component declaration and
port map of the \texttt{user\_logic} core inside
\texttt{sp101\_axi\_v1\_0\_S00\_AXI.vhd}.},label={lst:userlogic}]
-- Component declaration
component user_logic
    generic (
        A : natural := 10
    );
    port (
        run             : in  std_logic;
        reset           : in  std_logic;
        bus2mem_en      : in  std_logic;
        bus2mem_we      : in  std_logic;
        ck              : in  std_logic;
        bus2mem_addr    : in  std_logic_vector(A-1 downto 0);
        bus2mem_data_in : in  std_logic_vector(31 downto 0);
        sp2bus_data_out : out std_logic_vector(31 downto 0);
        done            : out std_logic
    );
end component;

-- Port map
U : user_logic
    generic map (
        A => 10
    )
    port map (
        run             => slv_reg0(0),
        reset           => slv_reg0(1),
        bus2mem_en      => slv_reg0(2),
        bus2mem_we      => slv_reg0(3),
        ck              => S_AXI_ACLK,
        bus2mem_addr    => slv_reg1(9 downto 0),
        bus2mem_data_in => slv_reg2,
        sp2bus_data_out => core_sp2bus_data_out,
        done            => core_done
    );
\end{lstlisting}

% =============================================================================
\section{Part 1 -- Behavioral Simulation (10 pts)}
% =============================================================================

\texttt{tb/tb\_user\_logic.vhd} drives the same \texttt{run},
\texttt{reset}, \texttt{bus2mem\_*} signals the assignment's xsim
script forces, looping over a single vector list
\texttt{SSQ\_X\_VALUES} and asserting that \(\mathrm{mem}[128]\)
holds \(x\cdot x\) after each run.  The vector list (Table~\ref{tab:vec})
includes the prompt example \((x=7)\), three corner cases \(\{0,1,2\}\),
small/medium positive values, and two large operands that stress the
32-bit datapath.

\begin{table}[H]
    \centering
    \caption{Test vectors driven by \texttt{tb\_user\_logic} (and by
             the on-board C app).  All execute the program
             \texttt{sc}~\(x\)~\texttt{ssq}~\texttt{halt} and verify
             \(\mathrm{mem}[128]=x^{2}\).}
    \label{tab:vec}
    \small
    \begin{tabular}{rrl}
        \toprule
        \(x\) & \(x^{2}\) & note \\
        \midrule
        \texttt{7}     & \texttt{49}         & prompt example \\
        \texttt{0}     & \texttt{0}          & zero edge case \\
        \texttt{1}     & \texttt{1}          & identity edge case \\
        \texttt{2}     & \texttt{4}          & small \\
        \texttt{3}     & \texttt{9}          & small \\
        \texttt{10}    & \texttt{100}        & round decade \\
        \texttt{15}    & \texttt{225}        & medium \\
        \texttt{16}    & \texttt{256}        & power of two \\
        \texttt{255}   & \texttt{65025}      & byte max \\
        \texttt{1000}  & \texttt{1000000}    & round thousand \\
        \texttt{32767} & \texttt{1073676289} & 15-bit max \\
        \texttt{46340} & \texttt{2147395600} & near \(\sqrt{2^{31}}\) \\
        \bottomrule
    \end{tabular}
\end{table}

The wave script
\href{run:../stackprocessor_101/scripts/sim_waves.tcl}{\texttt{scripts/sim\_waves.tcl}}
sets the data-path radices to \emph{unsigned} so the values read
directly as decimal (much easier to confirm \(7 \to 49\) than
\texttt{0x7}~\(\to\)~\texttt{0x31}).  Opcode signals like \texttt{ir}
stay hex because they are encoded constants, not numbers.

Figures~\ref{fig:sim1}--\ref{fig:sim3} show the full xsim run spanning
all twelve scenarios; the integer \texttt{checks} counter climbs
\(0\to 12\) while \texttt{errors} stays at~\(0\).  The
\texttt{ssq\_datapath} group in Figure~\ref{fig:sim3} is the clearest
proof of functional correctness: for every scenario the wave shows
\texttt{temp1}~=~\(x\), \texttt{mem\_data\_in}~=~\(x^{2}\) on the
\texttt{SSQ6} write edge, and \texttt{mem\_data\_out}~=~\(x^{2}\) on
the post-run read-back through the bus.  Figures~\ref{fig:simres1} and
\ref{fig:simres2} reproduce the Tcl-console transcript, listing the
twelve \texttt{[PASS]} notes with both decimal and hex values and the
\texttt{TESTBENCH PASSED} banner with \texttt{checks executed : 12,
errors : 0}.

\begin{figure}[H]
    \centering
    \optionalfig[\textwidth]{sim1.png}{xsim wave with bus signals and
        \texttt{checks} counter, scenarios 0--11.}
    \caption{xsim waveform, overview.  Top group shows
             \texttt{ck}, \texttt{run\_sig}, \texttt{reset\_sig},
             \texttt{bus2mem\_en/we}; the bus group shows
             \texttt{bus2mem\_addr}, \texttt{bus2mem\_data\_in} and
             \texttt{sp2bus\_data\_out} in decimal.  The
             \texttt{checks} integer (third pane from the top of the
             middle group) walks \(0,1,\dots,11\) across the run; the
             \texttt{errors} integer is constant at~\(0\).}
    \label{fig:sim1}
\end{figure}

\begin{figure}[H]
    \centering
    \optionalfig[\textwidth]{sim2.png}{xsim wave with processor
        registers and \texttt{ssq} datapath, scenarios 0--11.}
    \caption{xsim waveform, processor view.  \texttt{ir} (hex) walks
             through \texttt{SC \(\to\) SC2 \(\to\) \(\ldots\) \(\to\)
             SC5 \(\to\) SSQ \(\to\) SSQ2 \(\to\) \(\ldots\) \(\to\)
             SSQ6 \(\to\) HALT} every scenario.  \texttt{sp} bounces
             \(128\to 129\) once per scenario (\texttt{sc} bumps it,
             \texttt{ssq} leaves it).  \texttt{temp1} (decimal) holds
             the operand \(x\) and \texttt{mem\_data\_in} (decimal)
             carries \(x^{2}\) into the BRAM on the \texttt{SSQ6}
             write edge.}
    \label{fig:sim2}
\end{figure}

\begin{figure}[H]
    \centering
    \optionalfig[0.95\textwidth]{sim3.png}{Zoom on the
        \texttt{ssq\_datapath} and \texttt{self\_check} groups.}
    \caption{\texttt{ssq\_datapath} zoom (decimal radix).  Reading
             left to right the testbench drives \(x=0,7,1,2,3,10,15,
             16,255,1000,32767,46340\); each \texttt{temp1} sample is
             followed one cycle later by the matching
             \texttt{mem\_data\_in}~=~\(x^{2}\) and (after the
             read-back) \texttt{mem\_data\_out}~=~\(x^{2}\).  The
             \texttt{checks} counter climbs \(0\to 12\) and
             \texttt{errors} stays \(0\), proving every vector
             matches.}
    \label{fig:sim3}
\end{figure}

\begin{figure}[H]
    \centering
    \optionalfig[\textwidth]{sim_res1.png}{xsim Tcl console:
        per-scenario \texttt{[PASS]} notes.}
    \caption{xsim Tcl console output for all twelve scenarios.  Every
             vector reports
             \texttt{[PASS] ssq x=\(X\) -> \(X^{2}\) addr=128
             got=\(X^{2}\) expected=\(X^{2}\) (0x\dots)}.}
    \label{fig:simres1}
\end{figure}

\begin{figure}[H]
    \centering
    \optionalfig[0.9\textwidth]{sim_res2.png}{xsim Tcl console:
        final summary banner.}
    \caption{Final \texttt{tb\_user\_logic} banner.
             \texttt{checks executed : 12}, \texttt{errors : 0},
             \texttt{TESTBENCH PASSED -- all cases ok}.}
    \label{fig:simres2}
\end{figure}

% =============================================================================
\section{Part 3.2 -- Bare-metal C Application (3 pts)}
% =============================================================================

The bare-metal application
\href{run:../stackprocessor_101/sw/main.c}{\texttt{sw/main.c}} drives
the AXI4-Lite peripheral through the exact handshake described in the
PDF and runs the same twelve vectors as the simulator (so the on-board
UART log matches the simulator output line-for-line).  The key driver
helpers are reproduced in Listing~\ref{lst:cmain}; note the
``never assert \texttt{CTRL\_RESET} during \texttt{mem\_read}'' rule,
which is the failure mode that silently masked the dump in early
bring-up.

\begin{lstlisting}[caption={Core of \texttt{main.c}: bus write, bus
read, processor reset, and the per-scenario flow that proves
\(\mathrm{addr}\,128 = x^{2}\).},label={lst:cmain}]
/* slv_reg0 control bits */
#define CTRL_RUN        (1U << 0)
#define CTRL_RESET      (1U << 1)
#define CTRL_BUS_EN     (1U << 2)
#define CTRL_BUS_WE     (1U << 3)

/* IMPORTANT: never assert CTRL_RESET while reading.  Pulse reset only
 * in processor_reset(); use plain CTRL_BUS_EN (+CTRL_BUS_WE for
 * writes) for normal bus access. */
static void mem_write(u32 addr, u32 data)
{
    SP101_mWriteReg(BASEADDR, REG_CTRL,     CTRL_BUS_EN | CTRL_BUS_WE);
    SP101_mWriteReg(BASEADDR, REG_BUS_ADDR, addr & 0x3FFU);
    SP101_mWriteReg(BASEADDR, REG_BUS_DIN,  data);
}

static u32 mem_read(u32 addr)
{
    SP101_mWriteReg(BASEADDR, REG_CTRL,     CTRL_BUS_EN);
    SP101_mWriteReg(BASEADDR, REG_BUS_ADDR, addr & 0x3FFU);
    for (volatile int s = 0; s < 64; s++) { }
    return SP101_mReadReg(BASEADDR, REG_SP_DOUT);
}

static void processor_reset(void)
{
    SP101_mWriteReg(BASEADDR, REG_CTRL, CTRL_RESET | CTRL_BUS_EN);
    for (volatile int s = 0; s < 64; s++) { }
    SP101_mWriteReg(BASEADDR, REG_CTRL, CTRL_BUS_EN);
}

static void processor_run_and_wait(void)
{
    SP101_mWriteReg(BASEADDR, REG_CTRL, CTRL_RUN);
    while ((SP101_mReadReg(BASEADDR, REG_DONE) & 0x1U) == 0U) { }
    SP101_mWriteReg(BASEADDR, REG_CTRL, 0U);
}

/* Load `sc x  ssq  halt` at memory[0..3]. */
static void load_ssq_program(u32 x)
{
    const u32 prog[4] = { OP_SC, x, OP_SSQ, OP_HALT };
    for (u32 i = 0; i < 4U; i++) {
        mem_write(i, prog[i]);
    }
}

/* Per-scenario: reset -> poison mem[128] -> load program -> run ->
 * read mem[128] -> compare against x*x. */
static unsigned run_ssq_scenario(u32 x_in)
{
    const u32 expected = x_in * x_in;
    const u32 sentinel = 0xDEADBEEFU;

    processor_reset();
    mem_write(STACK_BASE, sentinel);   /* would survive a no-op bug */
    load_ssq_program(x_in);
    processor_run_and_wait();

    u32 got = mem_read(STACK_BASE);
    /* xil_printf("address 128 = %u  (expected %u = %u^2)  %s",
     *            got, expected, x_in, (got==expected)?"PASS":"FAIL"); */
    return (got == expected) ? 0U : 1U;
}
\end{lstlisting}

% =============================================================================
\section{Part 3.3 -- Hardware Results (UART, 4 pts)}
% =============================================================================

The same twelve scenarios were re-executed on the Cora~Z7-07S over the
AXI4-Lite interface (base \texttt{0x43C3\_0000}).  The UART transcript
(Figure~\ref{fig:vitis}, captured at \(115\,200\) Bd through
\texttt{ps7\_uart\_0}) reports every read-back of \(\mathrm{mem}[128]\)
as \texttt{PASS} -- decimal first, hex on the next line -- and ends
with \texttt{Summary : 12/12 scenarios PASSED}, matching the simulation
result line-for-line.  Listing~\ref{lst:uart} reproduces the captured
trace verbatim so the prompt's
``\(\mathrm{address}\,128\) contains the square'' requirement is
visible in text as well as in the screenshot.

\begin{figure}[H]
    \centering
    \optionalfig[0.45\textwidth]{vitis.png}{Serial-monitor capture of
        the twelve \texttt{ssq} scenarios on Cora~Z7-07S.}
    \caption{UART trace from the bare-metal \texttt{sp101 + ssq} test
             app on the Cora~Z7-07S.  Every scenario shows
             \texttt{address 128 = \(x^{2}\)} (decimal) and matching
             hex on the next line; the summary line at the bottom
             reads \texttt{Summary : 12/12 scenarios PASSED}.  Note
             the prompt case at the top: \texttt{address 128 = 49
             (expected 49 = 7\textasciicircum{}2) PASS}.}
    \label{fig:vitis}
\end{figure}

\begin{lstlisting}[language={},caption={Captured UART trace (verbatim
from the Cora~Z7-07S run; abbreviated -- full text in
Fig.~\ref{fig:vitis}).},label={lst:uart}]
==========================================================
   ECEC 661 FINAL - Stack Processor 101 + ssq
   IP base address : 0x43C30000
==========================================================

---- Scenario: ssq x=7 -> 49 ----
  program : sc 7  ssq  halt
    address 128 = 49  (expected 49 = 7^2)  PASS
                  hex: got=0x00000031  expect=0x00000031

---- Scenario: ssq x=0 -> 0 ----     ...   PASS
---- Scenario: ssq x=1 -> 1 ----     ...   PASS
---- Scenario: ssq x=2 -> 4 ----     ...   PASS
---- Scenario: ssq x=3 -> 9 ----     ...   PASS
---- Scenario: ssq x=10 -> 100 ----  ...   PASS
---- Scenario: ssq x=15 -> 225 ----  ...   PASS
---- Scenario: ssq x=16 -> 256 ----  ...   PASS
---- Scenario: ssq x=255 -> 65025 ----            ...   PASS
---- Scenario: ssq x=1000 -> 1000000 ----         ...   PASS
---- Scenario: ssq x=32767 -> 1073676289 ----     ...   PASS

---- Scenario: ssq x=46340 -> 2147395600 ----
  program : sc 46340  ssq  halt
    address 128 = 2147395600  (expected 2147395600 = 46340^2)  PASS
                  hex: got=0x7FFEA810  expect=0x7FFEA810

----------------------------------------------------------
   Summary : 12/12 scenarios PASSED
----------------------------------------------------------
\end{lstlisting}

% =============================================================================
\section{Conclusion}
% =============================================================================

The Final design adds a 6-state \texttt{ssq} micro-program
(Listing~\ref{lst:ssq}, opcode \texttt{0x00000041}) to the assignment
PDF's Stack Processor~101, packages the result as the AXI4-Lite IP
\texttt{sp101\_axi} (Listing~\ref{lst:userlogic} for the component
declaration and port map required by Part~2), and exercises twelve
vectors in both behavioural simulation
(Figures~\ref{fig:sim1}--\ref{fig:simres2}: \texttt{checks=12,
errors=0}) and on the Cora~Z7-07S
(Figure~\ref{fig:vitis} / Listing~\ref{lst:uart}:
\texttt{12/12 scenarios PASSED}).  For the prompt example
\texttt{sc 7 ssq halt}, \(\mathrm{mem}[128]=49\) on both simulation and
hardware and \(\mathit{sp}=129\) points at the next available address,
which is the exact behaviour required by
\href{run:../prob.md}{\texttt{Final/real/prob.md}}.  Each scored
deliverable maps to a concrete artifact under
\href{run:../stackprocessor_101/}{\texttt{Final/real/stackprocessor\_101/}}
and is reproduced step-by-step in
\href{run:../README.md}{\texttt{Final/real/README.md}}.

\end{document}
